\section{局限性}
\label{sec:limitations}

\textbf{已完成与未完成.} 本研究完成了大规模计算验证：1\,344 次受控多模型实验、环境因素析因与统计推断、跨模型序列归因、多尺度感受野分析、候选模式挖掘、细胞背景比较与证据整合。本研究**未完成**：直接序列扰动实验、编辑效率湿实验验证、独立生物学验证，以及在本批次中退化的真正留一细胞系（LOCO）重训。因此平台输出的是可检验假设，而非已验证机制。

\textbf{案例数据与设计.} 本文以 DeepCRISPR 数据作为案例研究，而非平台适用性的普遍证明。本批次中 \emph{all}（留一细胞系）划分与 \emph{single} 划分的 $R^{2}$ 逐位相同（$\max|R^{2}_{\text{all}}-R^{2}_{\text{single}}|=0$），说明该分支在网格中退化：本文因此不把按 held-out 细胞系分组的数字解释为跨细胞系泛化证据；真正的 LOCO 验证需要重新执行训练配置（本文未进行，以避免修改训练体系）。

\textbf{数值稳定性.} Linear Regression 在 192 次运行中有 56 次发散。平台通过隔离（$|R^{2}|>10$）而非删除处理，并在所有统计中以"排除发散后"的口径报告；该处理仍会影响线性模型在环境效应与归因比较中的权重，相关结论以中位数与非参数统计为主。

\textbf{环境编码与效应量级.} 表观遗传信息以逐位点二值形式进入模型，可能丢失连续信号与更高阶组合信息；环境中性结论限于当前编码方式，不能外推为生物学无关。此外，四个环境因子的效应量级均小于 0.01，即使证据等级较高（Tier 1/2）也不表示效应强度大。

\textbf{归因方法的边界.} 归因幅值不构成统计显著性；注意力权重不区分方向，本文仅将其作为支持性信息，且在本案例中其峰值位置与梯度/树归因不一致，故未用于候选模式提取。CNN 归因对卷积核大小敏感（$k=7$ 与 $k=3/5$ 的位置谱相关性仅 0.21--0.23），单一配置的归因不应作为独立结论。不同模型的 importance 在模型内归一化，其绝对量级不可跨模型类直接比较。

\textbf{候选模式与富集.} 候选模式以 4\,nt 为主且高度集中于 PAM 邻近区域，其统计支持部分来自该区域的整体重要性；通过富集的模式目前仅在单一细胞系（HeLa）中达到显著，属上下文依赖候选；跨模型比较仅限 CNN 三个核配置（缺少 Transformer IG 产物），seed 级支持无法计算。因此"跨模型收敛的生物学 motif"这一表述在本数据上不成立。

\textbf{统计口径.} 边级 bootstrap 基于逐样本重采样，与区组析因方差分析在依赖结构假设上不同，两类结论的差异已在 \S\ref{sec:res-environment} 说明；因子级置信区间由模型级均值 bootstrap 得到（$n=7$ 个模型），属于跨模型不确定性而非样本级不确定性。

\textbf{细胞背景与机制.} 细胞系间的效应差异可部分由样本量、标签分布与序列组成差异解释；本文不对其作机制性断言。

\textbf{工具范围.} 当前模型针对 SpCas9 与 NGG PAM；NAG/NGA 等非经典 PAM、其他 Cas 变体与碱基编辑系统均未纳入。平台当前的"设计"能力为候选优先级排序，不包含序列从头生成或约束满足式设计。
